\subsection{Algebra: geometric normality and the original tensor localizations}\label{algebra-proof-070-algebra-geometric-normality-and-the-original-tensor-localizations}

This editorial verification covers \texttt{algebra.tex} 46878--47091 at Stacks Project authority \texttt{a044\allowbreak{}46e5\allowbreak{}7ec1\allowbreak{}fbc2\allowbreak{}52a8\allowbreak{}71af\allowbreak{}cec7\allowbreak{}752f\allowbreak{}b280\allowbreak{}7b14}, SHA-256 \texttt{FA8B\allowbreak{}B92E\allowbreak{}58A4\allowbreak{}F78A\allowbreak{}2BD0\allowbreak{}1B3B\allowbreak{}6A4A\allowbreak{}87DE\allowbreak{}0A0D\allowbreak{}279F\allowbreak{}5DD9\allowbreak{}0641\allowbreak{}B574\allowbreak{}DD5F\allowbreak{}BFFF\allowbreak{}A4F3}. The \href{https://github.com/stacks/stacks-project/blob/a04446e57ec1fbc252a871afcec7752fb2807b14/algebra.tex\#L46878}{original source} remains the authority. Eight reports retain units MC-STK-ERR-0735--0737 with six existing operations. A designation receives one additional ``by''; the elliptic coordinated assumption is recorded as optional. Complete mathematical arguments remain linked editorial material, without replacing the translated source.

The source dependencies read here are 8217--8235, 8274--8309, 10062--10081, 10151--10194, 10220--10255, 10512--10542 and 44511--44524. Earlier complete proofs in groups 198, 204, 206, 268, 273, 275 and 282, and the completed smooth-normality and faithful-normality ascent/descent notes, are receiving dependencies, not new discoveries. No content of the four new literature-index hits was read.

\subsubsection{The four field-extension tests and the original inseparable tensor}\label{algebra-proof-070-the-four-field-extension-tests-and-the-original-inseparable-tensor}

Let \(k\) be the original field and \(A\) the original \(k\)-algebra. Normality means that every prime localization is an integrally closed domain. The zero algebra has no primes and is normal; every scalar extension of it is again zero, so all the stated tests hold for it. This case introduces no exclusion into the source theorem.

For an arbitrary field extension \(E/k\), the map \(A\to E\otimes_kA\) is faithfully flat. For every \(A\)-module \(M\), the actual comparison \[
 (E\otimes_kA)\otimes_AM\xrightarrow{\sim}E\otimes_kM,
 \quad (e\otimes a)\otimes m\longmapsto e\otimes am
\] has inverse \(e\otimes m\mapsto(e\otimes1)\otimes m\). Tensoring vector spaces over \(k\) is exact and reflects zero when \(E\ne0\); a basis of \(E\) identifies this functor with a nonempty direct sum of copies of the underlying \(M\). This also proves faithful flatness after any of the specified algebra base changes below.

The implications from all field extensions to finitely generated ones, from those to finite purely inseparable ones, and from all extensions to \(k^{\mathrm{perf}}\), follow by inclusion of the quantified classes. If \(k'/k\) is finite purely inseparable, its unique embedding into a chosen perfect closure is a map of the original \(k\)-extensions. Thus \[
 k'\otimes_kA\longrightarrow k^{\mathrm{perf}}\otimes_kA
 \cong k^{\mathrm{perf}}\otimes_{k'}(k'\otimes_kA)
\] is faithfully flat. The displayed associativity map sends \(u\otimes(v\otimes a)\) to \(uv\otimes a\), with inverse \(u\otimes a\mapsto u\otimes(1\otimes a)\). Normality of the target descends by the completed original normal-descent proof. In characteristic zero, the finite purely inseparable extensions and the perfect closure are the identity extension, so this step contains no positive-characteristic exponent assumption.

Every extension \(E/k\) is the directed union of its finitely generated intermediate fields: take the field generated by each finite set of elements, and use the union of two finite sets for a common upper index. The comparison \[
 \operatorname{colim}_i(k_i\otimes_kA)\xrightarrow{\sim}E\otimes_kA
\] sends a represented tensor to that same tensor. Each tensor and each tensor relation involve finitely many elements and hence occur at one sufficiently large field. This proves both surjectivity and injectivity and respects addition and multiplication. Normality passes to this colimit by the complete original prime-localization proof in group 206: contract a target prime, identify its local ring with the colimit of the normal local domains, descend each finite monic equation to a later stage, and transport the resulting equality \(a=bd\) in the stage ring. No map of entire fraction fields is required. Thus the finitely generated test implies the test for all fields.

It remains to derive the finitely generated test from the finite purely inseparable one. For the original finitely generated \(K/k\), the source field-diagram lemma supplies \[
 \begin{matrix}K&\longrightarrow&K'\\ \uparrow&&\uparrow\\ k&\longrightarrow&k',\end{matrix}
\] where \(k'/k\) and \(K'/K\) are finite purely inseparable, and \(K'/k'\) is separable. The exact construction is already proved in group 282. It permits \(K'=k'K\), but does not assert that \(k'\otimes_kK\) itself is a field. Retain the multiplication map \[
 \mu:k'\otimes_kK\longrightarrow k'K,\qquad
 \sum_i\lambda_i\otimes u_i\longmapsto\sum_i\lambda_iu_i.
\] In characteristic \(p>0\), choose the original finite exponent \(p^e\) with \((k')^{p^e}\subset k\). Its full Frobenius sum is \[
 \left(\sum_i\lambda_i\otimes u_i\right)^{p^e}
 =1\otimes\sum_i\lambda_i^{p^e}u_i^{p^e}.
\] All mixed multinomial coefficients vanish in the stated characteristic. If the original sum lies in \(\ker\mu\), the right-hand scalar is zero in the embedded original \(K\). Thus its \(p^e\)-th power is zero. Conversely every nilpotent maps to zero in the field. The image is a finite integral domain algebra over \(K\) containing \(K\) and \(k'\), hence is the compositum field: for a nonzero element its monic minimal equation over \(K\) has nonzero constant term and expresses its inverse in the image. Therefore \(\ker\mu\) is exactly the nilradical and \((k'\otimes_kK)_{\mathrm{red}}\cong k'K\). This records the quotient and both field inclusions; it does not discard the nilpotents of the original tensor. In characteristic zero the diagram uses the identity purely inseparable maps.

The finitely generated separable extension \(K'/k'\) is the fraction field of a smooth \(k'\)-domain \(B\). This is the source's smooth-model result 44511--44524: choose a finite-type domain with that fraction field and restrict to a nonzero smooth neighbourhood of its generic point. The complete generic smoothness argument is retained in \texttt{ALGEBRA\_\allowbreak{}FIELD\_\allowbreak{}FORMAL\_\allowbreak{}SMOOTH\_\allowbreak{}NOTES\_\allowbreak{}20260929.\allowbreak{}md}.

The finite purely inseparable test makes \(k'\otimes_kA\) normal. Base change of the actual smooth map makes \[
 k'\otimes_kA\longrightarrow B\otimes_{k'}(k'\otimes_kA)
\] smooth, so the target is normal by the preceding smooth-ascent theorem. The map \(b\otimes(\lambda\otimes a)\mapsto b\lambda\otimes a\), with inverse \(b\otimes a\mapsto b\otimes(1\otimes a)\), identifies it with \(B\otimes_kA\). Localizing at the images of all nonzero original \(b\in B\) gives \[
 K'\otimes_{k'}(k'\otimes_kA)\cong K'\otimes_kA.
\] On fractions, \((b/s)\otimes a\) corresponds to \((b\otimes a)/(s\otimes1)\); the inverse preserves the same \(s\). Normality is preserved under this localization, including a zero result if it occurs. Finally \(K\otimes_kA\to K'\otimes_kA\) is the base change of the field extension \(K'/K\), so it is faithfully flat and normality descends. This proves the last nontrivial implication among the original four conditions.

\subsubsection{Localization and separable field extensions}\label{algebra-proof-070-localization-and-separable-field-extensions}

For any multiplicative subset \(U\subset A\), any field extension \(E/k\) has the exact comparison \[
 E\otimes_kU^{-1}A\xrightarrow{\sim}(1\otimes U)^{-1}(E\otimes_kA),
 \quad e\otimes(a/u)\longmapsto(e\otimes a)/(1\otimes u).
\] Its inverse sends \((\sum_i e_i\otimes a_i)/(1\otimes u)\) to \(\sum_i e_i\otimes(a_i/u)\). Tensor relations and the same multiplicative-set equality witnesses prove both inverse identities. This also handles \(0\in U\), when both sides are zero. Since a prime localization of a localized ring is the corresponding original prime localization, normality localizes. Applying the four-condition definition proves the original localization assertion for geometric normality.

Let \(K/k\) be the source's separable field extension, without a finite-generation assumption. Group 273 proves the original reducedness assertion, so \(T=k^{\mathrm{perf}}\otimes_kK\) is reduced. The purely inseparable extension \(k^{\mathrm{perf}}/k\) is integral and radicial, and its field base change has a unique prime. One can check this with the original elements: for any \(z=\sum_i\lambda_i\otimes u_i\in T\), some finite exponent \(p^e\) puts every \(\lambda_i^{p^e}\) in \(k\), giving the same full Frobenius sum as above in the original \(K\). If this scalar is zero, \(z\) is nilpotent; if it is nonzero, \(z\) is a unit with inverse its \((p^e-1)\)-st power divided by that scalar. Thus every nonunit is nilpotent, and the nilradical is the only prime. Reducedness then makes every nonzero element a unit. The tensor is nonzero by faithful scalar extension, hence is a field. In characteristic zero it is just \(K\). The perfect-closure test proves geometric normality of the original \(K\).

\subsubsection{The original local reduction and normal finite subalgebras}\label{algebra-proof-070-the-original-local-reduction-and-normal-finite-subalgebras}

Assume now that \(A\) is geometrically normal over \(k\), and \(B\) is normal. If either algebra is zero, their tensor is zero and normal. Otherwise take an original prime \(\mathfrak r\subset A\otimes_kB\) and its contractions \(\mathfrak p,\mathfrak q\). Both factor maps to the target localization invert the elements outside these contractions. The map \[
 A_{\mathfrak p}\otimes_kB_{\mathfrak q}
 \longrightarrow (A\otimes_kB)_{\mathfrak r},\qquad
 (a/u)\otimes(b/v)\longmapsto (a\otimes b)/(u\otimes v)
\] therefore exists. The middle ring is the localization of \(A\otimes_kB\) at the multiplicative set generated by \(u\otimes1\), \(u\notin\mathfrak p\), and \(1\otimes v\), \(v\notin\mathfrak q\). The inverse to this first localization comparison sends these two types of inverse generators to \((1/u)\otimes1\) and \(1\otimes(1/v)\), respectively. Further localization at the prime induced by \(\mathfrak r\) yields exactly \((A\otimes_kB)_{\mathfrak r}\), with both composites fixing every original fraction. This proves the source reduction with its actual contractions. The additional ``by'' at 47000 only repairs their designation.

The localized \(A_{\mathfrak p}\) is geometrically normal by the previous result and is a normal domain by the identity field test; \(B_{\mathfrak q}\) is a normal domain. Retain the source's names \(A,B\) for these reduced cases and its fraction fields \(K=\operatorname{Frac}(A)\), \(L=\operatorname{Frac}(B)\). Unit 0736 corrects only ``fractions fields'' at both occurrences.

For every finite set \(E\subset B\), let \(D_E=k[E]\) and take its integral closure \(B_E\) in its own fraction field \(\operatorname{Frac}(D_E)\subset L\). Since fields are Nagata and the earlier finite-type theorem applies, \(B_E\) is finite over \(D_E\), hence a finite-type normal \(k\)-domain. Each element of \(B_E\) is integral over \(B\) and belongs to \(L\), so normality of the unchanged domain \(B\) puts it in \(B\). If \(E\subset E'\), the fraction-field inclusion sends any element integral over \(D_E\) to an element integral over \(D_{E'}\), giving the actual inclusion \(B_E\subset B_{E'}\). This is directed and its union is \(B\), since it contains every finite set \(E\). Consequently \[
 A\otimes_kB\cong\operatorname{colim}_E(A\otimes_kB_E)
\] by the same finite-tensor and finite-relation argument. Group 206 then reduces the desired normality to the case of the original finite-type normal domain \(B\). No Noetherian assumption on \(A\) is introduced. Unit 0737 closes all three occurrences of ``sub algebra'', including the locus omitted from one report's proposed count.

\subsubsection{Both tensor localizations inside the total quotient ring}\label{algebra-proof-070-both-tensor-localizations-inside-the-total-quotient-ring}

Write \(C=A\otimes_kB\), with the original domains and \(B\) now finite type over \(k\). It is nonzero: tensoring two nonzero vector spaces over a field does not kill \(1\otimes1\). It is flat over either factor. In particular every nonzero \(a\in A\) gives a nonzerodivisor \(a\otimes1\) on \(C\), and every nonzero \(b\in B\) gives a nonzerodivisor \(1\otimes b\). Let \(U\) and \(V\) be these respective multiplicative subsets. Keep the exact identifications \[
 U^{-1}C\cong K\otimes_kB=:D,\qquad
 V^{-1}C\cong A\otimes_kL=:E,\qquad
 (UV)^{-1}C\cong K\otimes_kL=:F.
\] They take the fractions \((a\otimes b)/(u\otimes1)\), \((a\otimes b)/(1\otimes v)\), and \((a\otimes b)/(u\otimes v)\) to \((a/u)\otimes b\), \(a\otimes(b/v)\), and \((a/u)\otimes(b/v)\), respectively; their inverses keep the displayed denominators. Every denominator is a nonzerodivisor in \(C\), so all four rings inject into the actual \(Q(C)\). Thus the comparison direction is \[
 C\subset D,E\subset F\subset Q(C).
\] No map from all of \(Q(C)\) into \(F\) is assumed. Such a map fixing \(C\) need not exist: for \(A=k[x],B=k[y]\), \(x+y\) is invertible in \(k(x,y)=Q(C)\), but is not a unit in the factor-denominator localization \(k(x)\otimes_k k(y)\). Indeed a relation \((x+y)h=f(x)g(y)\) with nonzero one-variable denominators would, after setting \(y=-x\), give \(0=f(x)g(-x)\) in the domain \(k[x]\), a contradiction. This verifies the direction of the actual embeddings used in the proof.

Reducedness of \(C\), needed by the source's finite-minimal-prime criterion, is available before invoking that criterion: geometric normality of \(A\) makes \(E=A\otimes_kL\) normal and therefore reduced. The injection \(C\hookrightarrow E\) makes \(C\) reduced. The normal ring \(E\) here uses the lowercase ground field, exactly the correction retained as MC-STK-ERR-0735. In general \(A\) is not even a \(K\)-algebra, so the printed uppercase tensor base has no justification.

The ring \(D=K\otimes_kB\) is a finite-type algebra over the field \(K\), hence Noetherian and has finitely many minimal primes. Every minimal prime of \(C\) contracts to zero in the domain \(A\): flat going down would otherwise put a strictly smaller prime below it. Hence every such prime avoids \(U\). Localization gives a bijection between minimal primes of \(C\) and minimal primes of \(D\). For the converse direction, contract a minimal prime of \(D\), take a minimal prime of \(C\) below it, and localize; minimality forces equality, and contraction recovers the original prime. Thus \(C\) has finitely many minimal primes, with no assumed Noetherianity of \(C\).

Here is the total-quotient comparison needed to use normality of both localizations. For any multiplicative subset \(W\) of nonzerodivisors of a ring \(C\), multiplication by a nonzerodivisor of \(C\) remains injective after localization. Conversely, if \(b/w\) is a nonzerodivisor on \(W^{-1}C\), then \(b\) is a nonzerodivisor on \(C\): from \(bc=0\) one gets \((b/w)(c/1)=0\), then \(c/1=0\), and injectivity of \(C\to W^{-1}C\) gives \(c=0\). Therefore the mutually inverse total-quotient maps are \[
 Q(C)\longrightarrow Q(W^{-1}C),\qquad
 a/b\longmapsto(a/1)/(b/1),
\] \[
 Q(W^{-1}C)\longrightarrow Q(C),\qquad
 (a/u)/(b/v)\longmapsto av/(ub).
\] In the second formula, \(u,v\in W\) and the just-proved criterion ensures that the original \(b\) is a nonzerodivisor; thus \(ub\) is an allowed denominator. Localization equality witnesses verify well-definedness, and the fraction formulas verify both inverse identities. Apply this to \(W=U\) and \(W=V\). These are identifications inside the original \(Q(C)\), rather than replacements of its objects by an unproved ambient field.

\subsubsection{Normality of the first localization and the exact intersection}\label{algebra-proof-070-normality-of-the-first-localization-and-the-exact-intersection}

Geometric normality of \(A\) implies geometric normality of its localization \(K\), hence geometric reducedness of \(K\). Write \(K=\bigcup_iK_i\) over all finitely generated intermediate fields. Each \(K_i\) is geometrically reduced: for every field extension of \(k\), its scalar extension injects into that of \(K\), because tensoring over the field preserves the original inclusion. A subring of a reduced ring is reduced. Thus \(K_i/k\) is separable in the source's equivalent field criterion.

For each original \(K_i\), choose the smooth \(k\)-domain \(D_i\) with fraction field \(K_i\) supplied by 44511--44524. Then \(B\to D_i\otimes_kB\) is the actual base change of a smooth map and is smooth. Since the original \(B\) is normal, smooth ascent makes this tensor normal. Localizing at the images of \(D_i\setminus\{0\}\) gives the normal ring \(K_i\otimes_kB\), with the map \((d/s)\otimes b\leftrightarrow(d\otimes b)/(s\otimes1)\) and its exact inverse. The natural transition maps are induced by the original field inclusions; no chosen smooth models need have compatible transition maps. Their tensor colimit is the original \(D=K\otimes_kB\). Group 206 proves it normal. This argument uses the earlier smooth-normality proof, which is independent of the present tensor-normality theorem, so there is no circular application.

Both \(D\) and \(E\) are now normal. Prove their intersection in the actual common ring \(F\) by tensoring the original vector-space sequence \[
 0\longrightarrow A\longrightarrow K\longrightarrow K/A\longrightarrow0
\] first with \(B\), then with \(L\). The comparison induced by \(B\hookrightarrow L\) is injective on \((K/A)\otimes_kB\), as tensoring over \(k\) is exact. If \(z\in K\otimes_kB\) has image in \(A\otimes_kL\), its image in \((K/A)\otimes_kL\) is zero. Injectivity forces its image in \((K/A)\otimes_kB\) to be zero, so exactness in the first row puts \(z\) in \(A\otimes_kB\). Conversely this last tensor lies in both rings. Hence, with all inclusions already proved, \[
 (K\otimes_kB)\cap(A\otimes_kL)=A\otimes_kB
\] inside \(F\), and therefore also inside \(Q(C)\).

Let the original \(z\in Q(C)\) be integral over \(C\). Its unchanged monic equation has coefficients in \(D\) as well as in \(E\). The exact total-quotient comparisons make it an element of both \(Q(D)\) and \(Q(E)\). A normal ring is integrally closed in its total quotient ring, by the complete original ideal and two-localization-multiplier proof in group 198. Therefore \(z\in D\) and \(z\in E\), and the proved intersection gives \(z\in C\).

We have proved that the original \(C\) is reduced, has finitely many minimal primes, and is integrally closed in its total quotient ring. The source criterion 8274--8309, with its complete product comparisons in group 204, therefore makes it normal. Its mechanism retains the actual coordinate idempotents of \(Q(C)=\prod_i\operatorname{Frac}(C/\mathfrak p_i)\): each satisfies \(T^2-T=0\) and so belongs to \(C\). The maps \(c\mapsto(e_ic)_i\) and \((c_i)_i\mapsto\sum_i c_i\) are inverse, giving \(C=\prod_i e_iC\). Each factor embeds in its original field and is integrally closed there, as a monic equation in that factor can be placed in the corresponding coordinate of the product with all other lower coefficients zero. Thus the factors are normal domains. Returning through the proved colimit of finite normal \(B_E\), and then through the original prime localizations, proves the full stated theorem for arbitrary \(A,B\) with the original hypotheses.

\subsubsection{Both scalar actions in the separable algebraic comparison}\label{algebra-proof-070-both-scalar-actions-in-the-separable-algebraic-comparison}

Let \(k'/k\) be the original separable algebraic extension and \(A\) a \(k'\)-algebra. For a finite purely inseparable \(L/k\), the ring \(L'=k'\otimes_kL\) is a field: separability gives reducedness and the purely inseparable field extension gives a single prime, as proved above. The exact maps \[
 A\otimes_kL\xrightarrow{\sim}A\otimes_{k'}L',
 \quad a\otimes\ell\longmapsto a\otimes(1\otimes\ell),
\] \[
 a\otimes(r\otimes\ell)\longmapsto ar\otimes\ell
\] are inverse by tensor balancing. Thus geometric normality over \(k'\) gives normality of every such \(A\otimes_kL\), and the original four-condition criterion gives geometric normality over \(k\).

Conversely assume geometric normality over \(k\), and take the original field extension \(K/k'\). Set \(F_0=k'\otimes_kk'\), acting on \(K\otimes_kA\) by \[
 (r\otimes s)(u\otimes a)=ur\otimes sa.
\] Let \(F_0\to k'\) be the original multiplication map. Then \[
 (K\otimes_kA)\otimes_{F_0}k'
 \xrightarrow{\sim}K\otimes_{k'}A,
 \quad (u\otimes a)\otimes c\longmapsto uc\otimes a,
\] with inverse \(u\otimes_{k'}a\mapsto(u\otimes a)\otimes1\). For well-definedness, multiplying on the left by \(r\otimes s\) gives \(urc\otimes sa=urcs\otimes a\), equal to multiplication of \(c\) by \(rs\) on the right. Conversely \((r\otimes1)-(1\otimes r)\) has zero image under multiplication, so the inverse respects the \(k'\)-balanced relation. These calculations verify both composites and both original scalar actions.

The earlier diagonal-localization proof, group 275, identifies \(F_0\to k'\) as a localization. In the finite separable case, retain a primitive element \(\alpha\), its monic minimal polynomial \(P(T)=(T-\alpha)Q(T)\), and its nonzero value \(Q(\alpha)\). The corrected original polynomial map is \[
 k'[T]/(P)\longrightarrow k'\otimes_kk',\qquad
 \sum_i a_iT^i\longmapsto\sum_i a_i\otimes\alpha^i.
\] It is an isomorphism by the basis \(1,\alpha,\ldots,\alpha^{d-1}\), retaining the summation fixed earlier. Inverting \(Q\) forces \(T=\alpha\), since \((T-\alpha)Q=0\), and gives exactly evaluation into \(k'\); the inverse is the coefficient inclusion. The full denominator rule is \(f(T)/Q(T)^n\mapsto f(\alpha)/Q(\alpha)^n\). For infinite separable algebraic \(k'/k\), take its finite intermediate fields and the multiplicative closure of all these original finite-stage sets in \(F_0\). Each element maps to a nonzero scalar. Every kernel element lies in a finite tensor stage and is killed there by one of the finite-stage denominators. Thus the localization map is injective after evaluation and surjective from the coefficient inclusion, proving the same comparison without a finite-generation assumption on \(k'/k\).

Base-changing this localization through the displayed \(F_0\)-action proves that \(K\otimes_{k'}A\) is a localization of the normal ring \(K\otimes_kA\). Hence it is normal. The original arbitrary \(K/k'\) proves geometric normality over \(k'\).

The algebraicity hypothesis cannot simply be removed from this equivalence. In characteristic \(p>0\), take the perfect field \(k=\mathbf F_p\), the rational fields \(k'=k(t)\), \(A=k(s)\), and the actual inclusion \(t\mapsto s^p\). The extension \(k'/k\) is separable, and \(A/k\) is separable, so \(A\) is geometrically normal over \(k\). But \[
 A\otimes_{k'}A\cong A[T]/(T^p-s^p)
                    =A[T]/((T-s)^p).
\] The first isomorphism sends the second factor's original \(s\) to \(T\) and the first factor's \(s\) to the coefficient \(s\). The polynomial \(X^p-t\) is irreducible over \(k(t)\), for instance by the Eisenstein argument at \(t\) in \(k[t]\); hence its degree and the indicated tensor basis are \(p\). The nonzero element \([T]-s\), nonzero in that monic basis, has \(p\)-th power zero. All intervening binomial coefficients vanish in characteristic \(p\). Thus this scalar extension is not reduced and not normal, proving that \(A\) is not geometrically normal over \(k'\). This verifies the source hypothesis and the need for its actual diagonal-localization argument; it is not a proposed extension or a new source correction.

\subsubsection{Report accounting and propagation}\label{algebra-proof-070-report-accounting-and-propagation}

OCC-01027 receives the minimal designation copyedit at 47000. OCC-01028 and OCC-12447 retain both operations of 0736. OCC-01029 and OCC-12448 retain all three operations of 0737. OCC-01031 and OCC-12446 retain the mathematically necessary tensor-base correction 0735. OCC-01030 is optional: the coordinated sentence already supplies the finite-type assumption through an ordinary elliptical predicate; adding ``is'' is readable but does not repair a missing mathematical hypothesis or an unintelligible sentence. It is not integrated as another required fix.

The source's omitted receiving comparisons have been proved with their actual rings, scalar actions, nonzero conditions, denominators, nilpotents and prime maps. The finite purely inseparable diagram receives group 282; reducedness and separability receive groups 268 and 273; the infinite diagonal localization receives group 275; the tensor-normality proof receives groups 198, 204 and 206, finite-type Nagata stability, smooth normal ascent and faithful normal descent. All full proofs remain editorial. No novelty or whole-chapter completion claim is made. Review continues at 47092.
